\newpage
\section{DESCRIPCIÓN GENERAL}


\subsection{Contexto del Problema}

El mercado secundario de boletería digital para espectáculos de entretenimiento y deportes opera actualmente
sobre profundas carencias 
de información y fallos de seguridad que afectan de forma sistemática a los consumidores, organizadores y
artistas. 
Aunque la transición tecnológica desde las entradas impresas tradicionales hacia los boletos con código QR
dinámico o estático 
optimizó los canales de distribución primaria, también introdujo nuevos vectores de explotación fraudulenta
que las bases de datos 
relacionales centralizadas y las pasarelas de pago tradicionales no han logrado erradicar por construcción.

A nivel global, la industria de la reventa de boletos es un sector con una valoración estimada de entre
$3\,000$ y más de $30\,000$\,millones de dólares anuales \cite{ticketfairy2024}. Paralelo a este crecimiento,
los esquemas de fraude transaccional y falsificación de entradas generan pérdidas globales que superan los
$1\,600$\,millones de dólares al año, afectando de manera directa a cerca de $5$\,millones de usuarios que
adquieren boletos fraudulentos en canales no oficiales \cite{dataintelo2024}. Esta problemática se sustenta
principalmente en tres puntos de fricción técnica y operativa:
\begin{enumerate}
    \item \textbf{Duplicación y Reventa Concurrente:} Un código QR estático o una representación en PDF de un
boleto constituye una cadena de datos binaria susceptible de ser clonada indefinidamente a costo marginal
cero. Si un revendedor comercializa el mismo archivo a múltiples compradores de forma paralela en mercados
informales, el sistema de validación centralizado de la ticketera local solo detectará la duplicidad en el
acceso físico al evento, denegando el ingreso al resto de los compradores de buena fe.
    \item \textbf{Asimetría de Liquidación Financiera:} Las pasarelas de pago Web2 tradicionales liquidan los
fondos de transacciones internacionales en ventanas temporales que oscilan entre $72$ y $120$ horas hábiles,
mientras que el boleto digital es transferido al comprador de forma inmediata. Esta brecha temporal permite a
revendedores maliciosos vender entradas, recibir el capital y posteriormente realizar contracargos bancarios
o reclamaciones de fraude, dejando al comprador final sin el activo y sin los fondos.
    \item \textbf{Ausencia de Trazabilidad Pública de la Cadena de Propiedad:} Al no existir un registro
público y transparente que certifique el historial de transferencias del boleto desde su acuñación primaria
por el organizador hasta el consumidor final, resulta imposible verificar la legitimidad de un boleto
adquirido en el mercado secundario antes de la validación en puerta.
\end{enumerate}

A nivel regional y local en Colombia, la debilidad técnica de las plataformas tradicionales ha propiciado
escenarios de fraude y 
cartelización de alta visibilidad. Un caso representativo ocurrió en el año 2020, cuando la Superintendencia
de Industria y 
Comercio (SIC) sancionó a la Federación Colombiana de Fútbol (FCF) y a la tiquetera oficial Ticketshop con
una multa superior a 
los $4{,}6$\,millones de dólares por la estructuración de una red de reventa masiva y desvío de más de
$42\,000$ entradas para las 
eliminatorias de la Copa Mundial de la FIFA Rusia 2018, reportando sobrecostos especulativos de hasta el 350\,\% sobre la tarifa 
nominal \cite{sic2020}. De igual forma, en conciertos multitudinarios celebrados en el país durante el año
2022, las autoridades 
locales reportaron miles de casos de estafas masivas por duplicidad de códigos QR distribuidos a través de
plataformas informales de 
reventa P2P, evidenciando que las tiqueteras convencionales no disponen de mecanismos para invalidar copias
previas al instante de una 
transferencia.

\subsection{Formulación del Problema}

El problema de investigación que aborda este trabajo se formula bajo la siguiente interrogante: 
\textit{¿cómo diseñar e implementar un middleware B2B de arquitectura Web2.5 que se acople a las plataformas
de boletería 
digital existentes para garantizar, por construcción y de forma atómica, la unicidad y autenticidad del
boleto en el mercado 
secundario, permitiendo regular de manera flexible la especulación y asegurar la liquidación inmediata de
pagos y comisiones, 
sin forzar a los usuarios a interactuar directamente con la complejidad técnica de la Web3?}

Este planteamiento da a proponer tres condiciones de diseño: la inhabilitación física de códigos anteriores
debe ocurrir de 
forma atómica con la confirmación de pago; el sistema debe operar como un servicio de bajo costo que se
acople a la infraestructura 
existente en lugar de reemplazarla; y la propiedad del ticket debe validarse on-chain con el fin de eliminar
la necesidad de un 
intermediario de custodia financiera.

\subsection{Propuesta de Solución}

La propuesta consiste en una plataforma de boletería digital híbrida bajo una arquitectura Web 2.5
estructurada como un middleware B2B.
El diseño del sistema no busca competir de manera directa con las ticketeras consolidadas del mercado, lo
cual generaría resistencia 
comercial e incrementaría la fricción técnica para el usuario final que suele tener desconfianza o
desconocimiento frente a los 
criptoactivos y billeteras Web3. En su lugar, el prototipo se concibe como un servicio auxiliar de bajo costo
transaccional con 
el que las ticketeras existentes se acoplan mediante APIs REST. 

Las plataformas de tickets tradicionales preservan su estructura Web2 (catálogo, reserva de asientos físicos,
frontend tradicional 
y autenticación JWT). Sin embargo, delegan en el middleware la liquidación transaccional y la transferencia
de propiedad on-chain 
sobre la red Stellar. Cada boleto se modela como un NFT en el entorno de Soroban \cite{sdf2024soroban}
identificado por la 
tupla \((\texttt{ticket\_root\_id}, \texttt{version})\). Durante la compra en reventa, el middleware
construye un XDR que ejecuta 
de forma atómica el patrón \textit{burn/remint} (invalidando la versión $v$ y emitiendo la versión $v+1$ al
comprador) y la 
distribución del pago en tokens estables USDC \cite{sdf2024stablecoins}.

La especulación de precios se maneja como una política moldeable. En lugar de imponer reglas que limiten la
flexibilidad del 
negocio, la restricción de precio máximo de reventa y el porcentaje de comisiones se establecen como opciones
configurables 
en el contrato inteligente al momento de su inicialización, supeditadas al acuerdo comercial exclusivo o a
las políticas 
operativas de la ticketera organizadora del evento.

\subsection{Justificación de la Solución}

La selección Stellar como blockchain y su motor de contratos inteligentes \textit{Soroban} como
infraestructura de soporte se justifica 
por sus propiedades orientadas al rápido procesamiento de transacciones financieras de manera económica:
\begin{itemize}
    \item \textbf{Latencia de Finalidad Determinista:} El algoritmo de consenso SCP de Stellar
\cite{mazieres2015stellar} 
    garantiza la finalidad absoluta de una transacción en un rango de $3$ a $5$ segundos. Esta velocidad es
coherente 
    con los tiempos de respuesta exigidos en pasarelas de pago y evita demoras en los accesos físicos del
aforo.

    \item \textbf{Bajo Costo Transaccional:} La comisión base en Stellar es de $100$ stroops
($\approx$\,0{,}02 COP a 
    la tasa de referencia), lo que permite procesar miles de transacciones de reventa y distribución de
regalías sin afectar
    el valor comercial de la entrada.

    \item \textbf{Atomicidad Nátiva y USDC via SAC:} A través del Stellar Asset Contract (SAC)
\cite{sdf2024sep41}, las 
    operaciones financieras de pago en USDC y las operaciones de estado del boleto (quema de la versión
anterior y 
    re-minteo del nuevo NFT) se ejecutan dentro del mismo marco de invocación transaccional de Soroban. Si
alguna de las 
    transferencias o validaciones de propiedad falla, la transacción se revierte en su totalidad, eliminando
la necesidad 
    de construir sistemas complejos de depósito en garantía (\textit{escrow}) que son vulnerables a fallos y
retención 
    indebida de fondos.

    \item \textbf{Arquitectura de Proyección Relacional:} El acoplamiento Web2.5 se apoya en un indexador
asíncrono que proyecta
     el historial on-chain en una base de datos relacional PostgreSQL. Esto permite que las aplicaciones web
de las tiqueteras 
     realicen consultas con latencias de milisegundos, ofreciendo una experiencia convencional de alta
velocidad de lectura sin 
     obligar al usuario a lidiar directamente con las llamadas de solo lectura a la blockchain.

    \item \textbf{Alta Capacidad de Procesamiento (TPS):} El protocolo Stellar soporta una tasa nominal de
procesamiento de 
    entre $1\,000$ y $3\,000$ transacciones por segundo (TPS) según la complejidad de las operaciones
\cite{binance2024stellar}. 
    Aunque esta cifra es inferior a la capacidad teórica de redes centralizadas tradicionales como Visa, que
reporta picos de 
    hasta $65\,000$\,TPS con promedios de procesamiento reales de entre $1\,700$ y $4\,000$\,TPS
\cite{visa2024tps}, 
    resulta órdenes de magnitud superior a otras redes blockchain públicas como Ethereum 1.0 (que oscila
entre $15$ y $30$\,TPS) 
    y es suficiente para absorber las demandas de alta concurrencia transaccional en eventos de boletería
masivos.

\end{itemize}

\subsection{Descripción del Proyecto}

\subsubsection{Objetivo General}

Desarrollar una solución B2B de reventa de tickets, basada en la tecnología de contratos inteligentes de
Stellar, que garantice la unicidad y autenticidad de cada activo digital, reduciendo al mínimo el fraude y la
especulación en la reventa de entradas para eventos de entretenimiento.

\subsubsection{Objetivos Específicos}

\begin{itemize}\setlength{\itemsep}{2pt}
    \item Crear un sistema de reventa B2B de tickets basado en contratos inteligentes y NFTs que aseguren
autenticidad, 
    originalidad y trazabilidad.
    \item Desarrollar un contrato inteligente que genere tickets como NFTs, permita la transferencia de
propiedad y 
    automatice la distribución de comisiones.
    \item Implementar un prototipo que simule el proceso completo de reventa en la red de prueba de Stellar,
desde la 
    venta hasta la compra.
    \item Gestionar el funcionamiento de la herramienta a través de una implementación web que permita la
simulación y 
    visualización del modelo y de su integración con la cadena de bloques.
\end{itemize}

\subsubsection{Entregables, Estándares y Justificación}

La Tabla~\ref{tab:entregables} consolida los entregables del trabajo de grado, los estándares aplicados como
marco de referencia 
y la justificación de cada uno.

\begin{table}[H]
\centering
\small
\renewcommand{\arraystretch}{1.3} 
\definecolor{rowgray}{gray}{0.96}
\caption{Entregables y estándares aplicados.}
\label{tab:entregables}
\begin{tabularx}{\textwidth}{>{\raggedright\arraybackslash}p{4.2cm} >{\raggedright\arraybackslash}p{4.2cm}
>{\raggedright\arraybackslash}X}
\toprule
\rowcolor{rowgray}
\textbf{Entregable} & \textbf{Estándar de referencia} & \textbf{Justificación} \\
\midrule
Memoria de Trabajo de Grado (este documento) & Formato institucional PUJ & Documento autocontenido que
sintetiza el trabajo realizado. \\
\rowcolor{rowgray}
Plan de Gestión del Proyecto (SPMP) & Plantilla institucional & Trazabilidad de la planeación, hitos, riesgos
y presupuesto. \\
Especificación de Requisitos (SRS) & Plantilla institucional & Definición detallada de RF, RNF, casos de uso
y criterios de aceptación. \\
\rowcolor{rowgray}
Especificación del Diseño (SDD) y Propuesta de TG & Modelo C4 + secuencias UML & Vistas de contexto,
contenedores, componentes y despliegue. \\
Suite de tres contratos Soroban + 58 pruebas & Rust + \texttt{soroban-sdk} 23 & Materialización del modelo
NFT con versionado, \textit{burn/remint} y NFT SEP-41. \\
\rowcolor{rowgray}
Backend Node.js + indexador & Express + Prisma + PostgreSQL & API REST, XDR proxy \textit{stateless} y
Sincronización on-chain/off-chain. \\
Frontend web (React/Vite) & SPA responsive + Freighter & Flujo E2E para usuarios finales y operadores. \\
\rowcolor{rowgray}
Integración continua & GitHub Actions & Compilación y pruebas automatizadas en cada \textit{push}. \\
\bottomrule
\end{tabularx}
\end{table}

Los artefactos de gestión, requisitos, diseño y verificación se organizan de manera reproducible y auditable
a partir del código 
fuente público y de la documentación entregada como anexos.
